\documentclass[11pt,a4paper]{article}
\usepackage[utf8]{inputenc}
\usepackage[T1]{fontenc}
\usepackage{lmodern}
\usepackage[ngerman]{babel}
\usepackage{amsmath,amssymb,amsthm}
\usepackage{geometry}
\geometry{a4paper, margin=2.5cm}
\usepackage{hyperref}
\usepackage{xcolor}
\usepackage{booktabs}
\usepackage{array}

\definecolor{navyblue}{RGB}{20, 45, 85}
\definecolor{accentblue}{RGB}{0, 102, 204}

\hypersetup{
  colorlinks=true,
  linkcolor=navyblue,
  citecolor=navyblue,
  urlcolor=accentblue,
  pdftitle={Klaerung der Kugeln},
  pdfauthor={Thomas Hoffbauer}
}

\newtheorem{theorem}{Satz}[section]
\newtheorem{lemma}[theorem]{Lemma}
\newtheorem{definition}[theorem]{Definition}
\newtheorem{corollary}[theorem]{Korollar}
\newtheorem{remark}[theorem]{Bemerkung}

\title{\textbf{\Large Kl\"arung der Kugeln}\\[0.45em]
  \large Freeze von Gate~3 und Er\"offnung des Fortsetzungsprojekts\\
  zur geometrischen Kl\"arung der Bamberg-Kugel}
\author{\textbf{Thomas Hoffbauer}\\[0.2em]
  \small Bamberg}
\date{9. Oktober 2026}

\begin{document}

\maketitle

\begin{abstract}
Gate~3A--3E der klassisch-analytischen $E_8$-Br\"ucke ist eingefroren.
Der kritische Pfad ist in Lean~4.34.1 gegen Mathlib derselben Fassung
kernelverifiziert. Sein Axiom-Footprint besteht aus \texttt{propext},
\texttt{Classical.choice} und \texttt{Quot.sound}. Die Riemannsche
Vermutung wird dabei nicht bewiesen, sondern exakt rekodiert.
Dieses Dokument setzt das fr\"uhere Manuskript \emph{Riemann Final}
au\ss er Kraft, soweit jenes einen Beweis der Vermutung beansprucht,
und er\"offnet das Fortsetzungsprojekt \emph{Kl\"arung der Kugeln}.
Die Firewall lautet
\[
  \text{Gate 3} \;\not\Rightarrow\; \text{BK-I2.7D}
  \;\not\Rightarrow_{\text{derzeit}}\; \text{BK-I2.7E}.
\]
Die Zahlen $240$, $12$ und $18$ bleiben getrennt.
\end{abstract}

\section{Gegenstand und Abgrenzung}

Die Kette, die mit diesem Freeze festgeschrieben wird, reicht von der
Gittergeometrie bis zur Nullstellengeometrie:
\begin{align*}
  \mathbb{Z}^{8}\text{-Poisson}
  &\longrightarrow E_8\text{-Poisson}\\
  &\longrightarrow \Theta_{E_8}\in M_4(\mathrm{SL}_2(\mathbb{Z}))\\
  &\longrightarrow \Theta_{E_8}=E_4\\
  &\longrightarrow r_{E_8}(2n)=240\,\sigma_3(n)\\
  &\longrightarrow D_{E_8}(s)=240\,\zeta(s)\zeta(s-3)\\
  &\longrightarrow \Lambda_{E_8}(4-s)=\Lambda_{E_8}(s)\\
  &\longrightarrow \text{Gamma-Br\"ucke}\\
  &\longrightarrow \text{Nullstellengeometrie}.
\end{align*}
Sie ist eine Rekodierung der klassischen Theorie in zwei um drei
verschobene Zeta-Familien. Sie ist keine Evidenz f\"ur die Wahrheit
der Riemannschen Vermutung und kein Satz \"uber dreidimensionale
Kugelpackungen.

\begin{remark}[Was nicht mitgemeint ist]
  \emph{Riemann Final} vom 8.~Oktober 2026 formuliert einen spektralen
  Beweis der Vermutung. Dieser Anspruch geh\"ort nicht zum
  zertifizierten Pfad. Ma\ss geblich ist der \"Aquivalenzsatz von
  Gate~3E. Zwei isolierte \texttt{sorry}-Deklarationen,
  \texttt{integer\_lattice\_poisson\_8} und die rohe
  Plancherel-Fassung in \texttt{E8PQC.lean}, liegen au\ss erhalb dieses
  Pfads und werden durch den Freeze weder geschlossen noch als Beleg
  verwendet.
\end{remark}

\section{Der eingefrorene Block}

Die folgenden Aussagen stehen als Theoreme in den Dateien
\begin{quote}\raggedright
\texttt{E8ThetaQExpansion.lean},\ 
\texttt{E8Dirichlet.lean},\ 
\texttt{E8DirichletFunctionalEquation.lean},\ 
\texttt{E8DirichletGammaBridge.lean}\ und
\texttt{E8DirichletZeroGeometry.lean}.
\end{quote}
$\Theta_{E_8}$ bezeichnet die Theta-Reihe auf der oberen Halbebene,
$E_4$ die normierte Eisensteinreihe vom Gewicht~$4$,
$\Lambda$ Mathlibs \texttt{completedRiemannZeta} und
$\Gamma_{\mathbb{R}}$ Delignes reellen Gamma-Faktor.

\begin{theorem}[Gate 3A und 3B]
  F\"ur jede positive ganze Zahl $n$ gilt
  $r_{E_8}(2n)=240\,\sigma_3(n)$.
  F\"ur $\operatorname{Re}s>4$ konvergiert die Dirichlet-Reihe absolut,
  und
  \[
    \sum_{n=1}^{\infty} r_{E_8}(2n)\,n^{-s}
    = 240\,\zeta(s)\zeta(s-3).
  \]
  Die linke Seite ist der Wert der totalisierten Funktion
  $D_{E_8}(s)=240\,\zeta(s)\zeta(s-3)$.
\end{theorem}

\begin{theorem}[Gate 3C]
  Die vervollst\"andigte Funktion
  $\Lambda_{E_8}(s)=240\,\Lambda(s)\Lambda(s-3)$
  erf\"ullt $\Lambda_{E_8}(4-s)=\Lambda_{E_8}(s)$ f\"ur jedes
  $s\in\mathbb{C}$.
  $D_{E_8}$ ist \"uberall meromorph. L\"angs der punktierten Umgebung
  von $s=1$ gilt
  \[
    \lim_{\substack{s\to 1\\ s\neq 1}} D_{E_8}(s)
    = 240\,\zeta'(-2),
  \]
  und $(s-4)D_{E_8}(s)$ strebt punktiert gegen $240\,\zeta(4)\neq 0$.
\end{theorem}

\begin{theorem}[Gate 3D]
  Sei $G(s)=\Gamma_{\mathbb{R}}(s)\Gamma_{\mathbb{R}}(s-3)$.
  Aus $G(s)\neq 0$ folgt $\Lambda_{E_8}(s)=G(s)\,D_{E_8}(s)$.
  Aus $G(s)\neq 0$ und $G(4-s)\neq 0$ folgen die multiplikative
  Gleichung $G(s)D_{E_8}(s)=G(4-s)D_{E_8}(4-s)$ und
  \[
    D_{E_8}(4-s)
    = \frac{G(s)}{G(4-s)}\,D_{E_8}(s).
  \]
  Die Nullstellen der totalisierten Funktion $G$ sind genau die Stellen
  $s=-2n$ und $s=3-2n$. Jede solche Stelle ist eine
  Unstetigkeitsstelle. Die Gleichung ohne diese
  Nichtverschwindensbedingung ist falsch; der Beleg ist $s=-2$.
\end{theorem}

\begin{definition}[Nichttriviale $E_8$-Nullstelle]
  $E_8^{\mathrm{nt}}(s)$ bedeutet: entweder
  $\zeta(s)=0$ und $0<\operatorname{Re}s<1$, oder
  $\zeta(s-3)=0$ und $0<\operatorname{Re}(s-3)<1$.
\end{definition}

\begin{theorem}[Gate 3E]
  \label{thm:3e}
  F\"ur die totalisierte Funktion gilt f\"ur jedes $s\in\mathbb{C}$
  \[
    D_{E_8}(s)=0
    \iff
    \zeta(s)=0 \lor \zeta(s-3)=0.
  \]
  Insbesondere ist der Funktionswert $D_{E_8}(1)=0$.
  Aus $E_8^{\mathrm{nt}}(s)$ folgt $\operatorname{Re}s\neq 2$, ohne
  die Riemannsche Vermutung.
  Mathlibs Aussage \texttt{RiemannHypothesis} ist \"aquivalent zu
  \[
    \forall s,\ 
    E_8^{\mathrm{nt}}(s)
    \rightarrow
    \bigl(\operatorname{Re}s=\tfrac12 \lor \operatorname{Re}s=\tfrac72\bigr).
  \]
  Die Identit\"at $4-\rho=3+(1-\rho)$ ist eine Gleichung in
  $\mathbb{C}$. Sie transportiert keine Nullstellen von $D_{E_8}$.
\end{theorem}

\section{Drei ausgeschlossene Fehldeutungen}

\begin{enumerate}
  \item Totalisierung und meromorphe Lesart sind verschieden.
  Der Lean-Wert $D_{E_8}(1)$ ist $0$, weil $\zeta(-2)=0$.
  Der punktierte Grenzwert ist $240\,\zeta'(-2)$.
  Dass diese Ableitung nicht verschwindet, ist nicht behauptet.
  \item $\operatorname{Re}s=2$ ist die Symmetrieachse von
  $\Lambda_{E_8}$, nicht eine kritische Linie von $D_{E_8}$.
  Satz~\ref{thm:3e} schlie\ss t nichttriviale $E_8$-Nullstellen auf
  dieser Achse aus, ohne eine Nullstelle von $\zeta$ zu lokalisieren.
  \item Der st\"arkste Satz \"uber die Vermutung ist eine
  \"Aquivalenz. $E_8$ liefert zwei um drei verschobene Familien der
  klassischen Nullstellen, keine zus\"atzliche Evidenz.
\end{enumerate}

\section{Firewall zur Bamberg-Kugel}

Die offenen Gates des Kugelprogramms sind BK-I2.7D und BK-I2.7E.
BK-I2.7D verlangt eine diskrete Punktmenge in $\mathbb{R}^3$, zw\"olf
n\"achste Nachbarn und die Gram-Matrix der kuboktaedrischen
fcc-Schale. BK-I2.7E verlangt f\"ur ein zuvor definiertes
dreidimensionales Gitter $\Lambda_{\mathrm{BK},s}$ die Dichte
$\pi/\sqrt{18}$.

\begin{theorem}[Keine Implikation]
  Aus den S\"atzen von Gate~3 folgt weder BK-I2.7D noch BK-I2.7E.
  Der Koeffizient $r_{E_8}(2)=240$ ist die Kusszahl von $E_8$ in
  $\mathbb{R}^8$. Er erzwingt keine zw\"olf Nachbarn in
  $\mathbb{R}^3$. Die Zahl $18=54/3=180/3240$ z\"ahlt Richtungen in
  der endlichen Norm-$2$-Schale der $120$ Einheitsicosianen. Sie ist
  kein Kovolumen und keine Minimaldistanz.
\end{theorem}

\section{Er\"offnung: Kl\"arung der Kugeln}

Das Fortsetzungsprojekt beginnt bei BK-I2.7D. Der offene Claim besteht
aus drei Schritten in fester Reihenfolge.
\begin{enumerate}
  \item \textbf{Konstruktion.}
  Eine Menge $S$ und eine Abbildung $\Pi:S\to\mathbb{R}^3$,
  hingeschrieben aus den vier Rang-$3$-Sektoren.
  \item \textbf{Diskretheit und Kontakt.}
  Das Bild ist diskret, und seine Minimalschale besteht aus zw\"olf
  Punkten.
  \item \textbf{fcc-Identifikation.}
  Die Gram-Matrix des Bildes ist bis auf eine orthogonale
  Transformation und eine Umnummerierung die Gram-Matrix der
  fcc-Kontaktschale.
\end{enumerate}
Schritt~2 und Schritt~3 sind erst definiert, wenn $\Pi$ hingeschrieben
ist. Der Zentralkosinus benachbarter Ecken des regul\"aren Ikosaeders
ist $\sqrt{5}/5$. Die paarweisen Kosinus der fcc-Kontaktschale liegen
in $\{-1,-1/2,0,1/2\}$. Eine orthogonale Abbildung erh\"alt diese
Skalarprodukte und f\"uhrt die eine Schale nicht in die andere.
$\sqrt{5}/3$ ist der Betrag des Kosinus des Diederwinkels.

BK-I2.7E bleibt downstream. Die Kepler-Verbindung bleibt eine Hypothese.
Ein Gate~3F wird nicht er\"offnet. Solange die drei Schritte nicht als
eigene S\"atze vorliegen, wird aus Gate~3 keine Dichte, keine
Kepler-Aussage und keine physikalische Kugelinterpretation abgeleitet.
Gate~3 bleibt die unver\"anderte Referenzbibliothek.

\noindent
\begin{minipage}{\textwidth}
\textbf{Freeze-Status}

\vspace{0.7em}
\renewcommand{\arraystretch}{1.25}
\begin{tabular}{@{}ll>{\raggedright\arraybackslash}p{6.4cm}@{}}
  \toprule
  Komplex & Status & Charakterisierung \\
  \midrule
  Gate 3A--3E & eingefroren &
  $E_8$, Modulform, Dirichlet-Reihe, Gamma-Br\"ucke,
  Nullstellengeometrie. Kernelverifiziert. \\
  BK-I2.7D & Engpass &
  Zuerst $\Pi$, dann Diskretheit und zw\"olf Punkte, dann die
  fcc-Gram-Matrix. \\
  BK-I2.7E & downstream &
  Kovolumen und Dichte $\pi/\sqrt{18}$, abh\"angig von I2.7D. \\
  Kepler-Verbindung & Hypothese &
  Strukturelle Koinzidenz der Zahl $18$, kein Beweis. \\
  \bottomrule
\end{tabular}
\end{minipage}

\end{document}
